\documentclass[10pt,a4paper]{amsart}
\usepackage[margin=19mm]{geometry}
\usepackage{amsmath,amssymb,mathtools}
\usepackage[scr=rsfs,cal=euler]{mathalfa}
\usepackage{xfrac}
\usepackage[unicode,hidelinks,bookmarks=false]{hyperref}
\newcommand{\ie}{i.e.\ }
\newcommand{\cf}{cf.\ }
\newcommand{\resp}{respectively\ }
\newcommand{\Subsection}{\subsection}
\providecommand{\nobreakdash}{\nobreak\hbox{-}}
\setlength{\emergencystretch}{2em}
\multlinegap0pt
\usepackage{amssymb}
\newtheorem{theorem}{Theorem}
\newtheorem{proposition}{Proposition}
\renewcommand{\Re}{\operatorname{Re}}
\title{A Counterexample to a Problem of Pommerenke\\on Convex Functions in the Class \texorpdfstring{$\Sigma$}{Sigma}}
\author{Yuankai Guo, Xiaozhe Hu}
\date{Source: arXiv:2609.04279v1 (CC BY 4.0)}
\begin{document}
\maketitle

\noindent\emph{Source and licence.} A Counterexample to a Problem of Pommerenke\\on Convex Functions in the Class \texorpdfstring{$\Sigma$}{Sigma}. Version arXiv:2609.04279v1 (CC BY 4.0), distributed by arXiv under the Creative Commons Attribution 4.0 International licence, http://creativecommons.org/licenses/by/4.0/, as recorded in the arXiv licence metadata for this exact version and shown on the abstract page https://arxiv.org/abs/2609.04279v1. Full licence notice: \texttt{licenses/arxiv-2609.04279-CC-BY-4.0.txt}.\par\smallskip

\noindent\textit{Editorial scope.} Counted unit: the article's theorem thm:main (source0.tex lines 100--102). The article's complete original proof of it is delivered with it (source0.tex lines 261--303, one \texttt{proof} environment of 43 lines, which the article places away from the statement). No mathematics is written, repaired or re-derived: the statement and the proof are the article's own lines, in the article's own notation, and no retained line is substituted or re-flowed.  Retained uncounted as the source-local context the proof needs: lines 76--78; lines 160--166.  Nothing was omitted from any retained span, so the package carries no omission record.  No retained line contains a citation, so the wrapper carries no bibliography and no bibliography entry.  No retained line contains a figure environment, an image-inclusion command or drawing code, so no image is delivered and the package declares no asset dependency.  Editorial formatting only, in addition: an AMS \texttt{amsart} preamble that restates the article's own declarations and macros used by the retained lines, namely the environments \texttt{theorem}, \texttt{proposition} on separate counters, together with the standard packages those lines need.  The retained environments print in this document as Theorem 1, Proposition 1 in order of appearance, whereas the article prints the same items with the numbering of its own full document.  The complete native source of the article as distributed in the arXiv e-print for this exact version is bundled unchanged as \texttt{sources/209355/source0.tex}.
\begin{equation}\label{eq:convex-criterion}
\Re\left(1 + \frac{z f''(z)}{f'(z)}\right) > 0, \qquad \forall z \in \Delta^*.
\end{equation}

\begin{proposition}\label{prop:G}
Let $G(z)$ be defined via its derivative
\begin{equation}\label{eq:def-Gprime}
G'(z) = (1 - \beta z^{-3})^{2/3}, \qquad z \in \Delta^*,
\end{equation}
with the principal branch satisfying $G'(\infty) = 1$, and normalized so that $G(z) = z + O(z^{-2})$ as $z \to \infty$. Then $G$ is single-valued, belongs to $\Sigma$, and is convex in $\Sigma$.
\end{proposition}

\begin{theorem}\label{thm:main}
There exist convex functions $F, G \in \Sigma$ and a parameter $\lambda \in (0, 1)$ such that the convex linear combination $H = \lambda F + (1-\lambda)G$ is not convex.
\end{theorem}

\begin{proof}[Proof of Theorem~\ref{thm:main}]
We evaluate the convexity criterion at the test point
\begin{equation}
z_0 = \frac{1}{r} = \frac{400}{399} \in \Delta^*.
\end{equation}
At this point, $\beta z_0^{-3} = \beta r^3 = w$. By definition, $1 - w = q^3$, where $q = \frac{1}{3} - \frac{1}{6}i$. Since $\arg q \in (-\pi/6, 0)$, we have $3\arg q \in (-\pi/2, 0) \subset (-\pi, \pi)$. Under the principal branch specified in Proposition~\ref{prop:G},
\begin{equation}
G'(z_0) = (1 - w)^{2/3} = (q^3)^{2/3} = q^2 = \left(\frac{1}{3} - \frac{1}{6}i\right)^2 = \frac{1}{12} - \frac{1}{9}i.
\end{equation}
The first derivative of $H$ at $z_0$ is
\begin{align}
D := H'(z_0) &= \frac{3}{5} F'(z_0) + \frac{2}{5} G'(z_0) \nonumber\\
&= \frac{3}{5}(1 - \alpha r^2) + \frac{2}{5} q^2 \nonumber\\
&= \frac{3}{5}\left(1 - \left(\frac{8}{9} + \frac{4}{9}i\right)\left(\frac{399}{400}\right)^2\right) + \frac{2}{5}\left(\frac{1}{12} - \frac{1}{9}i\right) \nonumber\\
&= \frac{184794 - 557603i}{1800000} \neq 0.
\end{align}

For the second derivatives at $z_0$,
\begin{equation}
z_0 F''(z_0) = 2\alpha r^2, \qquad z_0 G''(z_0) = \frac{2\beta z_0^{-3}}{(1 - \beta z_0^{-3})^{1/3}} = \frac{2w}{(1 - w)^{1/3}} = \frac{2w}{q}.
\end{equation}
Setting
\begin{equation}
A := \frac{3}{5}\alpha r^2 + \frac{2}{5}\frac{w}{q} = \frac{3}{5}\alpha r^2 + \frac{2}{5}\left(\frac{1}{q} - q^2\right),
\end{equation}
we have $z_0 H''(z_0) = 2A$. The curvature quantity at $z_0$ can then be expressed as
\begin{equation}
1 + \frac{z_0 H''(z_0)}{H'(z_0)} = 1 + \frac{2A}{D} = \frac{D + 2A}{D},
\end{equation}
where
\begin{equation}
D + 2A = \frac{3}{5}(1 + \alpha r^2) + \frac{2}{5}\left(\frac{2}{q} - q^2\right).
\end{equation}
A direct calculation in $\mathbb{Q}(i)$ gives
\begin{equation}
\frac{D + 2A}{D} = \frac{-54160961609 + 690164496000\,i}{69013985609}.
\end{equation}
In particular,
\begin{equation}
\Re\left(1 + \frac{z_0 H''(z_0)}{H'(z_0)}\right) = -\frac{54160961609}{69013985609} \approx -0.784782405 < 0.
\end{equation}
Because this real part is strictly negative, $H$ violates the exterior convexity criterion \eqref{eq:convex-criterion}, which completes the proof.
\end{proof}

\end{document}
